\scanpage{224}

\mspage{98}

(prop. 4.3, cor. 4, page 36) qu'il en résulte que $(W,C,F)$ est un triple de Quillen clos. Cela achève la dém. de 7.1 et 7.2.

Résumons l'essentiel :

% slanted
\margin{\uline{NB} C'est \unsure{finalement|surtout} ce corollaire qui est le résultat principal de la présente section.}

\uline{Cor.} 7.4. Soit $M$ une catégorie satisfaisant les conditions a) b) c) d) e) (p. 76) et e$'$), f) ((7.2.1.) p. 88) enfin g) (7.2.6) i.e. stabilité de l'ens $\mathrm{Mon}$ des monomorphismes de $M$ par cochgt. de base. Soit $W$ \deleted{\unsure{dans}} $\subset\mathrm{Fl}\,M$ un localiseur dans $M$, qui est \uline{stable par} \uline{$\lim$} \uline{filtrantes} et de plus (condition plus subtile, mais \inserted{pratiquement} anodine) \uline{accessible}.

On pose

\[
C=\mathrm{Mono},\ \mathrm{TC}=C\cap W,\ \mathrm{TF}=C_*,\ F=(\mathrm{TC})_*.
\]

Alors : 1°) $(C,\mathrm{TF})$ est un couple de Quillen clos (énoncé indépendant de la donnée de $W$) et $=$ un couple de Quillen à factorisation fonctorielle.

% slanted
\margin{\uline{NB} Dans les cas 1° et 2°, $C,\mathrm{TC}=C\cap W$ sont Q-engendrés chacun par une \uline{petite} partie.}

2°) Pour \inserted{que $(\mathrm{TC},F)$ soit un couple} \deleted{$\mathrm{TC}$ est stable par cochangement de} de Quillen, il f. et s. que $\mathrm{TC}$ soit stable par cochangement de base. (Et alors il y a factorisation fonctorielle.)

3°) Pour que $(W,C,F)$ soit un triple de Quillen clos, il f. et s. que les conditions suivantes soient satisfaites

\[
\left\{\begin{aligned}
&\text{a) }\mathrm{TC}\overset{\mathrm{def}}{=}W\cap\mathrm{Mono}\quad\text{stable par cochgt. de base}\\
&\text{b) }\mathrm{TF}\overset{\mathrm{def}}{=}(\mathrm{Mono})_*=\text{ applications de \unsure{type} \unsure{injectif}}\\
&\qquad\subset W.
\end{aligned}\right.\tag*{(7.4.1)}
\]
% archive sheet number: 223
